% ============================================================
%  ch0_front  封面、内容提要、译者说明、原书页码对照表
% ============================================================

\begin{titlepage}
\centering
\vspace*{2.2cm}
{\Huge\bfseries\heiti Drinfeld chtoucas\par}
\vspace{0.3cm}
{\Huge\bfseries\heiti 与 Langlands 对应\par}
\vspace{0.9cm}
{\Large\heiti Tata 研究所讲义\par}
\vspace{1.4cm}
{\Large Laurent Lafforgue\par}
\vspace{0.3cm}
{\large Institut des Hautes Études Scientifiques\par}
\vspace{1.2cm}
{\large 全文汉译与译注\par}
\vspace{0.4cm}
{\large 数译\ \ 译注\par}
\vspace{1.6cm}
{\small
底本：Laurent Lafforgue，\emph{Cours à l'Institut Tata sur les chtoucas de Drinfeld et la correspondance de Langlands}，\\
Institut des Hautes Études Scientifiques，Bures-sur-Yvette，Juin 2002，\\
预印本编号 IHES/M/02/45，共 56 页。\par}
\vfill
{\small 二〇二六年十一月\par}
\end{titlepage}

\section*{内容提要}
\addcontentsline{toc}{section}{内容提要}

这是 Lafforgue 于 2001 年 2 月在孟买 Tata 基础科学研究所讲授的十二小时课程的法语笔记。课程只讲一件事：函数域上 Langlands 对应之证明的总体架构。全篇分六章。第一章给出对应本身的陈述、Drinfeld chtoucas 及其水平结构的定义、相应模叠的基本性质，并引入 Harder--Narasimhan 多边形截断与不动点计数公式。第二章说明截断后边界上同调的可忽略性，由此把上同调中的“本质部分”分离出来，并在其上得到 Hecke 代数作用的迹公式；这是全篇的技术核心。第三章用“完整同态”构造无水平结构的紧化，引入不分离的 fourre-tout 叠并分解其边界分层。第四章在 fourre-tout 内定义开紧化，陈述光滑性定理与主紧合性定理。第五章把紧合性化归为一个赋值判据，并证明 Hecke 对应下的全局稳定性定理。第六章用初等变换的序列造出满足判据的迭代 $\varphi$-格，从而补全第五章；其中修正了 \cite{lafforgue98} 末页由 Yakov Varshavsky 指出的一处小错。

原文是写给已经熟悉几何 Langlands 纲领的读者的，行文极简，许多论证只给“思路”而不给细节。译文中凡遇原文一笔带过、而读者可能在此卡住的地方，以脚注补出该步所依据的命题或直觉，并逐一注明来源层级。

\medskip
\noindent\textbf{关键词：}Langlands 对应；函数域；Drinfeld chtouca；Harder--Narasimhan 截断；紧化；Hecke 对应；迹公式。

\section*{译者说明}
\addcontentsline{toc}{section}{译者说明}

\textbf{底本与校对。}\ 译文以 IHES 预印本 IHES/M/02/45（2002 年 6 月，56 页）为唯一底本，逐句对译，未参照任何其他版本。底本 PDF 的文本层由 2002 年的 Ghostscript 以 Type3 位图字体生成，抽取后全是乱码，无法直接复制，故全部正文依页面图像逐页录入、核对。原书印刷页与本译本 PDF 页码相差 1：原书第 1 页是本译本 PDF 的第 2 页。每页起首处用以 \origpage{1} 这样的灰字标出原书页码，完整对照见下页《原书页码对照表》。

\textbf{编号系统。}\ 原文以罗马数字 I--VI 分章，章内不再分节；所有定理、命题、引理、推论、定义、注一律\textbf{不编号}，只用“Théorème.”“Proposition.”这样的字样起头。本译本原样保留这一体例：节号用罗马数字 I--VI，定理类环境也不编号。因此正文中不会出现“定理 3.2”式的引用；需要指称时，用“第 III 章的命题”这样的说法。公式同样不编号。

\textbf{术语。}\ 先查本库《数语规范》（\texttt{01\_corpus\_zh/数语规范/TERMS.tsv}），有定案者遵其定案；无定案者取中文数学文献的通行译法，首次出现时以括号附原文；尚无通行中译者保留原文，不作音译。全篇只有一种写法，术语总表见《工程文件/术语表.md》。以下几点需要事先说明：

\begin{itemize}
\item \emph{chtouca} 保留原文不译（本库 \texttt{SY02-18} 处理 \emph{shtuka} 时亦保留原文）。该词由 Drinfeld 所造，文献中另有 \emph{shtuka} 一种拼法，所指相同。
\item \emph{fourre-tout} 保留原文不译。这是 Lafforgue 自造的称呼，指一个把所有截断紧化都收作开子叠的不分离叠；译作“杂物袋”既失其名，又易与“栈”“叠”混淆。
\item \emph{cohomologie essentielle} 译“本质上同调”，\emph{cohomologie négligeable} 译“可忽略上同调”。这里的“可忽略”是相对于固定的秩 $r$ 而言的（见第 II 章定义），不是通常意义下的“零”。
\end{itemize}

\textbf{译注。}\ 只写三类：易错预警、含糊点拨、背景直觉。每条译注都以脚注给出，且必带 \textbf{[来源：……]}；凡属译者依上下文推断而无外部依据者，一律注明 \textbf{[译者按（上下文推断，无外部依据）]}。原文的笔误不擅自改动，照录并加注说明，同时登记在《工程文件/reports/02\_疑点清单.tsv》。

\textbf{文献。}\ 底本末尾只有一份九条的“简略书目”，照录不译，以便读者按原样检索。正文中的方括号引用（如 [Lafforgue, 1997]）一律改用 \LaTeX{} 的 \verb|\cite|，与书目中的条目一一对应。

\section*{原书页码对照表}
\addcontentsline{toc}{section}{原书页码对照表}

下表左列是原书印刷页码，中列是该页范围的内容，右列是本译本中与之对应的节。查原文时按下表定位即可。

\begin{longtable}{p{2.0cm}p{9.2cm}p{3.2cm}}
\toprule
原书页 & 该页范围的内容 & 本译本位置\\
\midrule
\endfirsthead
\toprule
原书页 & 该页范围的内容 & 本译本位置\\
\midrule
\endhead
\bottomrule
\endfoot
1--2 & 标题页、引言：课程的由来，六章内容提要 & 引言\\
3--11 & I.\ 证明的原理：Langlands 对应的陈述、chtoucas 的定义、模叠、截断、不动点计数 & 第 I 章\\
12--18 & II.\ 证明的大线条：边界可忽略性、本质上同调、迹公式、Hecke 代数的作用 & 第 II 章\\
19--26 & III.\ 无水平结构的紧化：完整同态、fourre-tout、边界分层 & 第 III 章\\
26--33 & IV.\ 紧化上的截断与水平结构：开紧化、光滑性定理、主紧合性定理 & 第 IV 章\\
34--41 & V.\ 紧合性的赋值判据；全局稳定化与局部稳定化 & 第 V 章\\
42--55 & VI.\ 赋值判据的证明：$\varphi$-空间、迭代 $\varphi$-格、退化子、初等变换序列 & 第 VI 章\\
56 & 简略书目（九条） & 参考文献\\
\end{longtable}

\newpage
\tableofcontents
\newpage
